\documentclass[11pt, a4paper]{ctexart}

\usepackage{assignment_style}

\begin{document}

% =========================================
% 1
% =========================================
\begin{problem}{1. 已知函数表如下：}
    \[
        \begin{array}{c|ccc}
            x_i & 101 & 102 & 103 \\ \hline
            f(x_i) & 10.04988 & 10.09950 & 10.14889
        \end{array}
    \]
    用三点公式求 $f'(101)$ 和 $f'(102)$。
\end{problem}

\begin{solution}
    节点等距，步长为
    \[
        h=1.
    \]
    三点前向公式为
    \[
        f'(x_0)\approx \frac{-3f(x_0)+4f(x_1)-f(x_2)}{2h},
    \]
    三点中心公式为
    \[
        f'(x_1)\approx \frac{f(x_2)-f(x_0)}{2h}.
    \]
    代入
    \[
        x_0=101,\qquad x_1=102,\qquad x_2=103
    \]
    及表中函数值，得
    \[
        f'(101)\approx \frac{-3\times 10.04988+4\times 10.09950-10.14889}{2}=0.049735,
    \]
    \[
        f'(102)\approx \frac{10.14889-10.04988}{2}=0.049505.
    \]
\end{solution}


% =========================================
% 3
% =========================================
\begin{problem}{3. 设 $f(x)=\dfrac{1}{(1+x)^2}$，取步长 $h=0.1$。}
    \begin{enumerate}[(1)]
        \item 用二点公式取 $x_0=0.5,\ x_1=0.6$ 计算 $f'(0.5)$，并估计误差；
        \item 用三点公式取 $x_0=0.4,\ x_1=0.5,\ x_2=0.6$ 计算 $f'(0.5)$，并估计误差。
    \end{enumerate}
\end{problem}

\begin{solution}
    已知
    \[
        f(x)=\frac{1}{(1+x)^2},
    \]
    故
    \[
        f'(x)=-\frac{2}{(1+x)^3},\qquad
        f''(x)=\frac{6}{(1+x)^4},\qquad
        f'''(x)=-\frac{24}{(1+x)^5}.
    \]

    \textbf{(1) 二点公式}

    由前向差分公式
    \[
        f'(x_0)\approx \frac{f(x_1)-f(x_0)}{h},
    \]
    取 $x_0=0.5,\ x_1=0.6$，得
    \[
        f'(0.5)\approx \frac{f(0.6)-f(0.5)}{0.1}
        =\frac{\dfrac{1}{1.6^2}-\dfrac{1}{1.5^2}}{0.1}
        =-0.53819.
    \]
    其截断误差满足
    \[
        R_1'(x_0)=\frac{h}{2}f''(\xi),\qquad \xi\in(0.5,0.6).
    \]
    在区间 $[0.5,0.6]$ 上，
    \[
        |f''(\xi)|\le \frac{6}{1.5^4},
    \]
    所以
    \[
        |R_1'(0.5)|\le \frac{0.1}{2}\cdot \frac{6}{1.5^4}=0.059259.
    \]

    \textbf{(2) 三点公式}

    由三点中心公式
    \[
        f'(x_1)\approx \frac{f(x_2)-f(x_0)}{2h},
    \]
    取 $x_0=0.4,\ x_1=0.5,\ x_2=0.6$，得
    \[
        f'(0.5)\approx \frac{f(0.6)-f(0.4)}{0.2}
        =-0.59790.
    \]
    截断误差为
    \[
        R_2'(x_1)=\frac{h^2}{6}f'''(\xi),\qquad \xi\in(0.4,0.6).
    \]
    又因为在区间 $[0.4,0.6]$ 上
    \[
        |f'''(\xi)|\le \frac{24}{1.4^5},
    \]
    所以
    \[
        |R_2'(0.5)|\le \frac{0.1^2}{6}\cdot \frac{24}{1.4^5}=0.007437.
    \]
\end{solution}


% =========================================
% 4(1)
% =========================================
\begin{problem}{4.(1) 用梯形公式和 Simpson 公式计算积分}
    \[
        I=\int_0^1 e^{-x}\,\dif x
    \]
    的近似值，并估计其截断误差。
\end{problem}

\begin{solution}
    对于区间 $[0,1]$，
    \[
        f(x)=e^{-x},\qquad a=0,\qquad b=1.
    \]

    \textbf{梯形公式}

    梯形公式为
    \[
        T(f)=\frac{b-a}{2}\bigl[f(a)+f(b)\bigr].
    \]
    因而
    \[
        T(f)=\frac{1}{2}(1+e^{-1})\approx 0.683940.
    \]
    误差项满足
    \[
        R_T=-\frac{(b-a)^3}{12}f''(\xi),\qquad \xi\in(0,1).
    \]
    由于
    \[
        f''(x)=e^{-x},\qquad \max_{[0,1]}|f''(x)|=1,
    \]
    故
    \[
        |R_T|\le \frac{1}{12}.
    \]

    \textbf{Simpson 公式}

    Simpson 公式为
    \[
        S(f)=\frac{b-a}{6}\left[f(a)+4f\left(\frac{a+b}{2}\right)+f(b)\right].
    \]
    所以
    \[
        S(f)=\frac{1}{6}\left(1+4e^{-1/2}+e^{-1}\right)\approx 0.632334.
    \]
    误差项满足
    \[
        R_S=-\frac{(b-a)^5}{2880}f^{(4)}(\xi),\qquad \xi\in(0,1).
    \]
    而
    \[
        f^{(4)}(x)=e^{-x},\qquad \max_{[0,1]}|f^{(4)}(x)|=1,
    \]
    从而
    \[
        |R_S|\le \frac{1}{2880}.
    \]
\end{solution}


% =========================================
% 5
% =========================================
\begin{problem}{5. 求 Newton-Cotes 系数 $C_k^{(3)}\ (k=0,1,2,3)$。}
\end{problem}

\begin{solution}
    四点闭型 Newton-Cotes 公式为
    \[
        \int_{x_0}^{x_3}f(x)\,\dif x \approx 3h\sum_{k=0}^{3}C_k^{(3)}f(x_k),
        \qquad x_k=x_0+kh.
    \]
    对应的系数可由 Lagrange 基函数积分得到，计算结果为
    \[
        C_0^{(3)}=\frac{1}{8},\qquad
        C_1^{(3)}=\frac{3}{8},\qquad
        C_2^{(3)}=\frac{3}{8},\qquad
        C_3^{(3)}=\frac{1}{8}.
    \]
    因而该公式即为 3/8 Simpson 公式：
    \[
        \int_{x_0}^{x_3}f(x)\,\dif x
        \approx \frac{3h}{8}\bigl[f(x_0)+3f(x_1)+3f(x_2)+f(x_3)\bigr].
    \]
\end{solution}


% =========================================
% 8
% =========================================
\begin{problem}{8. 求近似公式}
    \[
        \int_0^1 f(x)\,\dif x \approx \frac{1}{3}\left[2f\left(\frac{1}{4}\right)-f\left(\frac{1}{2}\right)+2f\left(\frac{3}{4}\right)\right]
    \]
    的代数精确度。
\end{problem}

\begin{solution}
    记右端为 $Q(f)$。依次验证单项式 $1,x,x^2,x^3,\ldots$。

    当 $f(x)=1$ 时，
    \[
        \int_0^1 1\,\dif x=1,\qquad
        Q(1)=\frac{1}{3}(2-1+2)=1.
    \]
    当 $f(x)=x$ 时，
    \[
        \int_0^1 x\,\dif x=\frac{1}{2},\qquad
        Q(x)=\frac{1}{3}\left(2\cdot \frac14-\frac12+2\cdot \frac34\right)=\frac{1}{2}.
    \]
    当 $f(x)=x^2$ 时，
    \[
        \int_0^1 x^2\,\dif x=\frac{1}{3},\qquad
        Q(x^2)=\frac{1}{3}\left(2\cdot \frac{1}{16}-\frac{1}{4}+2\cdot \frac{9}{16}\right)=\frac{1}{3}.
    \]
    当 $f(x)=x^3$ 时，
    \[
        \int_0^1 x^3\,\dif x=\frac{1}{4},\qquad
        Q(x^3)=\frac{1}{3}\left(2\cdot \frac{1}{64}-\frac{1}{8}+2\cdot \frac{27}{64}\right)=\frac{1}{4}.
    \]
    但当 $f(x)=x^4$ 时，
    \[
        \int_0^1 x^4\,\dif x=\frac{1}{5},
    \]
    而
    \[
        Q(x^4)=\frac{1}{3}\left(2\cdot \frac{1}{256}-\frac{1}{16}+2\cdot \frac{81}{256}\right)=\frac{37}{192}\neq \frac{1}{5}.
    \]
    因而该求积公式对所有不超过三次的多项式都精确，但对四次多项式不精确，所以其代数精确度为
    \[
        3.
    \]
\end{solution}


% =========================================
% 16(1)
% =========================================
\begin{problem}{16.(1) 确定下列公式中的求积系数，使其具有尽可能高的代数精确度：}
    \[
        \int_{-h}^{h}f(x)\,\dif x \approx A_1f\left(-\frac{h}{2}\right)+A_2f(0)+A_3f\left(\frac{h}{2}\right).
    \]
\end{problem}

\begin{solution}
    令公式对低次单项式尽可能精确。

    当 $f(x)=1$ 时，
    \[
        \int_{-h}^{h}1\,\dif x=2h=A_1+A_2+A_3.
    \]
    当 $f(x)=x$ 时，
    \[
        \int_{-h}^{h}x\,\dif x=0=-\frac{h}{2}A_1+\frac{h}{2}A_3.
    \]
    当 $f(x)=x^2$ 时，
    \[
        \int_{-h}^{h}x^2\,\dif x=\frac{2h^3}{3}=\frac{h^2}{4}A_1+\frac{h^2}{4}A_3.
    \]
    解此方程组得
    \[
        A_1=\frac{4h}{3},\qquad
        A_2=-\frac{2h}{3},\qquad
        A_3=\frac{4h}{3}.
    \]
    因而求积公式为
    \[
        \int_{-h}^{h}f(x)\,\dif x
        \approx
        \frac{4h}{3}f\left(-\frac{h}{2}\right)
        -\frac{2h}{3}f(0)
        +\frac{4h}{3}f\left(\frac{h}{2}\right).
    \]

    检验 $f(x)=x^3$ 时，
    \[
        \int_{-h}^{h}x^3\,\dif x=0,
    \]
    而右端也因对称性等于 $0$，故公式对三次多项式仍精确。
    但对 $f(x)=x^4$，
    \[
        \int_{-h}^{h}x^4\,\dif x=\frac{2h^5}{5},
    \]
    而求积值为
    \[
        \frac{4h}{3}\left(\frac{h}{2}\right)^4+\frac{4h}{3}\left(\frac{h}{2}\right)^4=\frac{h^5}{6}\neq \frac{2h^5}{5}.
    \]
    因而该公式具有三次代数精确度。
\end{solution}

\end{document}
